% Einheit 5 von 5 – Kurvenuntersuchung im Sachzusammenhang (bank/kurvenuntersuchung/e5.jsonl, zusammenbau v0.1)
%% TODO Zeitmarke Einheit 5: Marken-Zeile nicht lesbar
%% TODO Zweigzeile Teil 1: was hier gelernt wird (Satz in Schülersprache aus dem Einheitstitel) – Einheit 5
\einheitenkopf[e5]{Einheit 5 von 5 · Kurvenuntersuchung im Sachzusammenhang}
\zweigzeile{Abitur GK}
\setcounter{aufgabe}{33}

%% TODO Titel: Ich-kann-Satz fehlt in der Bank (Platzhalter: Kettenname) – Nr. 34
%% TODO Anweisung: ein Satz über der ersten Teilaufgabe fehlt in der Bank – Nr. 34
\begin{aufgabe}{Sachzusammenhang}
%% TODO \rechenplatz{n}: Schrittzahl des Grundfalls fehlt in der Bank (nur bei fester Schreibform, 2.3 b)
\begin{teile}
\teil Kreuze an, was die Frage mathematisch bedeutet, ohne zu rechnen: „Wann wächst die Pflanze am schnellsten?“ \\ \kreuz{größter Wert: Hochpunkt, Ränder vergleichen} \\ \kreuz{stärkste Änderung: Extremum von $f'$, Wendepunkt} \\ \kreuz{keine Änderung: $f'(t) = 0$}
\teil Die Abbildung zeigt die Temperatur $T$ (in °C) eines Tees $t$ Minuten nach dem Aufgießen. Beschreibe den Verlauf im Sachzusammenhang.

\begin{ksys}[xmin=0,xmax=40,ymin=0,ymax=100,xstep=10,ystep=20,xlabel=t in min,ylabel=T in °C,ablesen] \funktionab{20+70*exp(-0.1*\x)}{T}{0}{40} \end{ksys}
\teil Die Abbildung zeigt die Höhe $h$ (in m) eines Heißluftballons $t$ Minuten nach dem Start mit dem Wendepunkt $W$. Gib die Bedeutung des Wendepunkts für die Fahrt an. \hfill (Abitur 2022 GK)

\begin{ksys}[xmin=0,xmax=20,ymin=0,ymax=90,xstep=5,ystep=10,xlabel=t in min,ylabel=h in m,ablesen] \funktionab{-0.02*\x^3+0.6*\x^2}{h}{0}{20} \wendepunkt{10}{40}{W} \end{ksys}
\teil Die Abbildung zeigt die Höhe $h$ (in cm) einer Schneedecke $t$ Stunden nach Beginn des Tauwetters. Lies ab, wie hoch der Schnee zu dem Zeitpunkt ist, an dem er am schnellsten taut.

\begin{ksys}[xmin=0,xmax=20,ymin=0,ymax=50,xstep=5,ystep=5,xlabel=t in h,ylabel=h in cm,ablesen] \funktionab{0.01*\x^3-0.3*\x^2+45}{h}{0}{20} \end{ksys}
\teil Die Oberkante eines Gartenbogens wird durch $f(x) = -0{,}2x^2 + 1{,}6x + 0{,}4$ beschrieben (1 LE = 1 m). Berechne die Koordinaten des Hochpunkts und gib die Höhe des Bogens an. \hfill (FHR 2023)
\teil Die Wassermenge in einem Speicher wird durch $W(t) = -0{,}5t^3 + 6t^2 + 100$ beschrieben ($t$ in Stunden, $0 \le t \le 10$, $W(t)$ in m³). Weise nach, dass sie nach $8$ Stunden am größten ist, und gib diesen Wert an.
\teil Die Konzentration eines Medikaments im Blut wird durch $c(t) = 6t \cdot \mathrm{e}^{-0{,}2t}$ beschrieben ($t$ in Stunden, $c(t)$ in mg pro Liter); ohne Nachweis gilt $c''(t) = (0{,}24t - 2{,}4) \cdot \mathrm{e}^{-0{,}2t}$. Bestimme den Zeitpunkt, zu dem die Konzentration am stärksten abnimmt – die notwendige Bedingung genügt –, und gib Konzentration und Änderungsrate zu diesem Zeitpunkt an. \hfill (Abitur 2019 GK)
\teil Die Zahl der Gäste in einem Freizeitpark wird durch $B(t) = -t^3 + 6t^2 + 15t + 100$ beschrieben ($t$ in Stunden nach der Öffnung, $0 \le t \le 8$). Wann sind die meisten Gäste im Park und wie viele?
\teil Eine Rennbahn für Spielzeugautos ist in der Draufsicht symmetrisch zur $x$- und zur $y$-Achse. Ihr oberer Verlauf wird auf $[-3; 0]$ durch $g(x) = \frac{1}{9}x^3 + 0{,}5x^2$ beschrieben, mit waagerechter Tangente bei $x = -3$; oberer und unterer Verlauf sind links und rechts durch je einen Halbkreis verbunden (1 LE = 1 m). Zeige, dass die Bahn auf eine Platte von $3{,}2$ m $\times$ $9{,}2$ m passt. \hfill (Abitur 2026 GK)
\teil Die Höhen zweier Pflanzen $A$ und $B$ werden durch $f(t)$ und $h(t)$ beschrieben ($t$ in Wochen, Höhe in cm); bis $t_s$ ist $A$ höher als $B$. Gegeben ist ein Lösungsweg: $d(t) = f(t) - h(t)$; $d'(t) = 0$ hat für $0 < t < t_s$ nur die Lösung $t_1 \approx 4{,}2$; $d''(t_1) < 0$; $d(t_1) \approx 2{,}6$. Beschreibe die Bedeutung von $d(t)$ und deute die Zahlen $4{,}2$ und $2{,}6$ im Sachzusammenhang. \hfill (Abitur 2024 GK)
\end{teile}
\end{aufgabe}

%% TODO Titel: Ich-kann-Satz fehlt in der Bank (Platzhalter: Kettenname) – Nr. 35
%% TODO Anweisung: ein Satz über der ersten Teilaufgabe fehlt in der Bank – Nr. 35
\begin{aufgabe}{Funktionswert im Sachzusammenhang deuten}
\begin{teile}
\teil Die Flughöhe eines Drachens wird durch $h(t)$ beschrieben ($t$ in Minuten nach dem Start, $h(t)$ in m). Deute $h(5) = 32$ im Sachzusammenhang.
\end{teile}
\end{aufgabe}

%% TODO Titel: Ich-kann-Satz fehlt in der Bank (Platzhalter: Kettenname) – Nr. 36
%% TODO Anweisung: ein Satz über der ersten Teilaufgabe fehlt in der Bank – Nr. 36
\begin{aufgabe}{Mittleren Wert einer periodisch schwankenden Größe am Graphen ablesen}
\begin{teile}
\teil Die Abbildung zeigt die Temperatur $T$ (in °C) an einem Sommertag $t$ Stunden nach Mitternacht. Lies den mittleren Wert ab, um den die Temperatur schwankt.

\begin{ksys}[xmin=0,xmax=24,ymin=0,ymax=25,xstep=4,ystep=5,xlabel=t in h,ylabel=T in °C,ablesen] \funktionab{18+4*sin(15*(\x-9))}{T}{0}{24} \end{ksys}
\end{teile}
\end{aufgabe}

%% TODO Titel: Ich-kann-Satz fehlt in der Bank (Platzhalter: Kettenname) – Nr. 37
%% TODO Anweisung: ein Satz über der ersten Teilaufgabe fehlt in der Bank – Nr. 37
\begin{aufgabe}{Sachzusammenhang – Fehler finden, Begründen, Anwendung, Darstellungswechsel}
\begin{teile}
\teil Der Graph der Höhe $h$ eines Ballons hat bei $t = 10$ einen Wendepunkt und bei $t = 20$ seinen Hochpunkt ($t$ in min). Tom schreibt: „Ab $t = 10$ sinkt der Ballon.“ Finde den Fehler und stelle richtig.
\teil Begründe, warum der höchste Wasserstand in einem Beobachtungszeitraum am Anfang oder am Ende liegen kann, obwohl der Graph dort keinen Hochpunkt hat.
\teil Die Oberkante einer Rampe für Spielzeugautos wird für $0 \le x \le 15$ durch $f(x) = -\frac{1}{250}x^3 + 0{,}09x^2 + 2$ beschrieben (1 LE = 1 cm); an den Enden ist sie waagerecht. Der Anstiegswinkel darf nirgends größer als $30°$ sein. Bestimme den Punkt mit dem größten Anstieg und prüfe die Bedingung; die hinreichende Bedingung ist nicht verlangt. \hfill (Abitur 2017 GK)
\teil Die Kosten für die Herstellung von $x$ m³ einer Farbe werden durch $K(x) = x^3 - 9x^2 + 30x + 16$ beschrieben ($0 \le x \le 9$, $K(x)$ in Tausend Euro); die Abbildung zeigt den Graphen von $K$. Der Erlös ist $E(x) = 24x$. Zeichne den Graphen von $E$ ein und bestimme aus der Darstellung, bei welchen Verkaufsmengen Gewinn erzielt wird. \hfill (Abitur 2018 GK)

\begin{ksys}[xmin=0,xmax=9,ymin=0,ymax=300,xstep=1,ystep=50,xlabel=x in m³,ylabel=K,ablesen] \funktionab{\x^3-9*\x^2+30*\x+16}{K}{0}{9} \end{ksys}
\end{teile}
\end{aufgabe}
